\chapter{Related Work}\label{chap:rw}

The automation of trading card sorting is a multi-disciplinary problem intersecting mechanical engineering, robotics, and computer vision. While the market for \acp{TCG} has existed for decades, automated solutions have only recently become viable due to advancements in deep learning and affordable rapid prototyping tools. This chapter reviews existing commercial systems and relevant academic literature, categorizing them into mechanical handling solutions and algorithmic recognition strategies.

\section{Commercial Solutions}
The current market for automated sorting is dominated by high-cost, closed-source proprietary systems. These machines are typically designed for high-volume businesses rather than individual collectors.

\subsection{Hardware-Based Sorters}
\begin{itemize}
    \item \bb{\ac{TCG} Machines:} The \ii{PhyzBatch-9000} by \ii{\ac{TCG} Machines} \cite{tcgmachines} represents the industrial standard. It addresses a specific computer vision challenge: holographic glare. By utilizing specialized LED arrays and diffusers, the machine controls lighting conditions to detect \enquote{foil} (holographic) cards, which are prone to over-exposure in standard camera setups.
    \item \bb{CardBot:} \ii{CardBot} \cite{cardbot} differs in its handling mechanism. Unlike friction-based feeders, which can scratch the surface of cards, \ii{CardBot} uses a pneumatic suction system. This approach minimizes physical contact with the card surface, which is essential to prevent damage to the cards.
    \item \bb{Magic Sorter:} \ii{Magic Sorter} \cite{magicsorter} offers a more compact alternative, focusing specifically on \ii{Magic: The Gathering} cards. It integrates a feeder mechanism with a custom software suite, demonstrating the tight coupling between hardware and proprietary databases common in commercial tools.
\end{itemize}

\subsection{Software Ecosystems}
Apart from hardware, pure software solutions attempt to streamline the cataloging process. \ii{Card Dealer Pro} \cite{carddealerpro} provides a computer-vision-based inventory management system. While effective, such solutions rely on the user to manually present cards to a webcam or scanner, lacking the robotic automation required for bulk processing.

\section{Algorithmic Approaches to Recognition}
The core challenge in automating \ac{TCG} sorting is the reliable identification of card artwork among tens of thousands of similar classes. Research in this domain generally falls into two categories: \ac{OCR} and \ac{VSL}.

% \subsection{OCR and Text-Based Identification}
% A traditional approach involves reading the card's title and set code. A case study by \ii{Oodles Technologies} for \ii{Erie \ac{TCG}} \cite{erietcg} explored this using \ac{NER} and \ac{OCR}. 

% While theoretically sound, this approach struggles in practice due to the \enquote{artistic} nature of \acp{TCG}. Text is often stylized, obscured by artwork elements, or printed on holographic backgrounds that confuse standard \ac{OCR} engines. Consequently, text extraction is often computationally expensive and brittle compared to visual matching.

\subsection{Visual Similarity and Metric Learning}
Modern approaches favor deep learning techniques that treat the card as a holistic visual fingerprint rather than a text document.

\subsubsection{Siamese Networks}
Lowhur \cite{lowhuryugioh} demonstrated a highly effective system for \ii{Yu-Gi-Oh!} cards achieving 99\% accuracy. The methodology moves away from standard classification (which requires retraining for every new card set) to \ii{Metric Learning}. 

By utilizing Siamese Networks and triplet loss functions, the model learns an embedding space where images of the same card are clustered together. This allows the system to identifying cards by comparing their feature vectors against a database, a technique known as \ii{One-Shot Learning}.

\subsection{Object Detection and Robotics Integration}
Recent academic work has focused on the pipeline of locating objects (detection) and manipulating them (robotics).

\subsubsection{YOLO for Real-Time Detection}
The \ac{YOLO} architecture is the standard for real-time localization.
\begin{itemize}
    \item \bb{Robotic Integration:} Yu \cite{pokerrobot} validated the integration of vision and pneumatics in a self-playing poker robot. The study successfully combined \ac{YOLO}v5 for card localization with a suction-cup robotic arm. This provides a direct precedent for the hardware-software integration proposed in this thesis, confirming that suction mechanisms are viable for thin card stock.
    \item \bb{Pipeline Validation:} Freitter \cite{idcardverification} proposes a pipeline for ID card verification that mirrors the requirements of \ac{TCG} sorting: segmentation using \ac{YOLO}v8 followed by classification using \ac{ResNet}. This \enquote{two-stage} approach: first finding the object, then identifying it; is robust against background noise and variable lighting.
\end{itemize}

\subsubsection{The Data Scarcity Problem}
A major hurdle in training neural networks for TCGs is the lack of labeled training data for the millions of unique cards in circulation. Pua \cite{pokemonstanford} addressed this in a study on \ii{Pokémon} card detection. Instead of manually photographing thousands of cards, the authors developed a \bb{Synthetic Data Generation} pipeline. By projecting digital card scans onto varying backgrounds and applying simulated lighting, noise, and occlusions, they successfully trained a robust \ac{YOLO}v8 model. This methodology highlights that synthetic data is not only a time-saver but a necessity for achieving high robustness in \ac{TCG} recognition.

\section{Summary}
Current commercial solutions like \ii{CardBot} provide robust hardware but remain inaccessible due to cost and closed ecosystems. Conversely, academic research provides individual building blocks: \ac{YOLO} for detection \cite{pokerrobot}, synthetic data for training \cite{pokemonstanford}, and Siamese networks for identification \cite{lowhuryugioh}. But rarely integrates them into a complete, open-source device. This thesis aims to bridge this gap by combining these state-of-the-art algorithmic techniques with a low-cost, replicable hardware design.